\problemname{I Don't Miss Pennies}

\newcommand{\hmaxn}{3 \cdot 10^{5}}
\newcommand{\hmaxp}{3\,000}

In November 2025, the United States minted its last pennies (one-cent
coins).  Canada has not minted pennies since 2012.

Billions of pennies remain in circulation, but pennies are
expected to fade from circulation over time.  Stores are expected to
continue to price items to the cent, as credit card transactions are
processed to the cent.  However, for cash transactions, stores are
expected to round to the nearest nickel (five cents).  Specifically, if
the final digit of a total purchase ends in $3$, $4$, $8$ or $9$ cents, the
total will be rounded up; if it ends in $1$, $2$, $6$ or $7$ cents, it will be
rounded down. Transactions ending in $0$ or $5$ cents are not rounded.

You realize that this provides an opportunity.  If you pay cash and
rearrange and group your purchases appropriately, you may be able to pay slightly
less for everything you buy!  Given the prices of the individual
items you wish to buy in cents, determine the maximum amount you could save by paying
solely in cash and rearranging and grouping your purchases optimally, compared to the
non-rounded total price you would pay if you bought all of the items 
with a credit card.

\section*{Input}
The first line of input contains a single integer $N$ ($1 \le N \le \hmaxn$), the number of items you intend to purchase.

The next line contains $N$ space-separated integers $p_i$ ($1\le p_i\le \hmaxp$), the prices of the items in cents.

\section*{Output}
Print a single integer: the difference between the total credit card price without
rounding and the lowest total cost that can be obtained by paying cash
and grouping purchases optimally.

\section*{Explanation of Sample Input 1}

For the first sample, one optimal grouping of items is
\[\{78, 999, 350\}, \{59, 173, 2995, 350\}, \{882, 298\}, \{1096\}, \{497\}.\]
The prices of each group are, respectively, $1427, 3577, 1180, 1096, 497$ and thus the total difference is
\[(1427 + 3577 + 1180 + 1096 + 497) - (1425 + 3575 + 1180 + 1095 + 495) = 7.\] 
You save $7$ cents by optimally grouping the items into several cash transactions instead of paying for everything with a credit card.

